% !LOOM digest: arxiv:2401.00002
% !LOOM prefix: arxiv-2401.00002v1
% !LOOM extracted-from: arxiv:2401.00002v1
% !LOOM method: extract
% !LOOM proofs: verbatim
% !LOOM created: 2026-09-18
% !LOOM numbering: emulated
% !LOOM tool: weft 0.1.0.dev0

\section{The vanishing}\label{arxiv-2401.00002v1-sec-1}

\begin{theorem}[{\cite[Theorem 1.1]{arxiv:2401.00002}}]\label{arxiv-2401.00002v1-thm-1.1}
\label{arxiv-2401.00002v1-thm:vanish}
Let $C$ be a weighted stable curve of genus $g$ with weights summing to less than one. Then $H^1(C, \omega_C) = 0$.
\end{theorem}

\begin{proof}
The weight condition forces every component to be rational, so the dualising sheaf is a sum of line bundles of negative degree. The vanishing for each summand is the local computation of \cite[Theorem 1.1]{other}, and the compatibility of the two normalisations is the point of \cite{third}.
\end{proof}

\section{A consequence}\label{arxiv-2401.00002v1-sec-2}

\begin{corollary}[{\cite[Corollary 2.1]{arxiv:2401.00002}}]\label{arxiv-2401.00002v1-cor-2.1}
\uses{arxiv-2401.00002v1-thm-1.1}
\label{arxiv-2401.00002v1-cor:rigid}
Such a curve admits no infinitesimal automorphism.
\end{corollary}

\begin{proof}
Immediate from Theorem~\ref{arxiv-2401.00002v1-thm-1.1} and Serre duality.
\end{proof}
